%=============================================================================
\section{Introduction}
%=============================================================================

This document is \textbf{not} a textbook. Textbooks are great---rigorous, comprehensive, and excellent at putting you to sleep at 2 AM. This is a \textit{lecture note}, written by someone who learned control the hard way: by breaking robots first, then reading the theory to figure out why.

The goal is simple: give you enough theoretical foundation to \textbf{design, tune, and debug controllers} on real hardware. We will cover the math, but for every concept we will always ask: ``What is this actually for when the robot hits the field?''

\subsection*{Who Is This For?}

Undergraduate students in their early and middle years. Prerequisites: calculus and basic linear algebra (taken or currently taking). If you know what a derivative is and can multiply matrices without breaking down, you are ready.

\subsection*{What Is Special About This Note's Structure?}

Almost no control textbook on the market has a similar content arrangement.

We start with \textbf{Digital Signal Processing}---because before you can control anything, you need clean data from your sensors. Then we learn how to \textbf{describe systems} mathematically---a fancy way of saying ``write down what your motor actually does.'' Next comes \textbf{Classical Control Theory}, PID's home turf---the bread and butter of competition robotics. After that, \textbf{Discretization and Implementation} bridges the gap between elegant math and the harsh reality of a microcontroller running at a fixed clock rate. Finally, \textbf{Modern Control Theory} gives you the tools that turn a ``working'' robot into a ``good'' robot: LQR, MPC, and Kalman filters.

Let us begin.
